% !TeX program = lualatex
% !TeX encoding = utf8
% !TeX spellcheck = uk_UA
% !BIB program = biber
% !TeX root =../LabWork.tex


\abstract{%
\begin{enumerate}
    \item Експериментально підтвердити та проаналізувати справедливість закону Ома для різних ділянок електричного кола (повного кола, зовнішньої ділянки, неоднорідної ділянки), визначити зв'язок між струмом, напругою та опором.

    \item Визначити внутрішній опір $R_i$ реального джерела електричної енергії.

    \item Дослідити режим узгодження потужностей джерела та споживача.

    \item Класифікувати досліджуване джерело за величиною його внутрішнього опору (тобто визначити, чи наближається воно до ідеального джерела напруги чи до ідеального джерела струму).

    \item Визначити умови роботи джерела (режими: холостого ходу, короткого замикання, узгодженого навантаження) та розрахувати максимальний коефіцієнт корисної дії (ККД) для кожного з режимів, з'ясувавши, за якого режиму ККД є найбільшим.
\end{enumerate}
}
%======================================================================================================================%

\keywords{}

\apparatus{}

\chapter{Джерела електричної енергії та їх характеристики}
\makeworktitle

\section{Теоретичні відомості}


%\begin{Definition}
Електричним струмом називають упорядкований рух електричних зарядів в провідному середовищі під дією електричного поля.

Мірою інтенсивності руху зарядів є \emph{сила струму} $I$~--- кількість заряду, що проходить через поперечний переріз провідника за одиницю часу:
\begin{equation}\label{Idef}
	I = \frac{dq}{dt}.
\end{equation}
Для постійного струму ($I = \mathrm{const}$):
\begin{equation}\label{Idef1}
	I = \frac{q}{t}.
\end{equation}
%\end{Definition}

\begin{wrapfigure}{l}{0.4\linewidth}\centering
	\begin{tikzpicture}[thick, every circuit symbol/.style={thick},large circuit symbols,]
        \def\h{1}
		\draw (0,-\h) coordinate (START) to [battery={info={$\mathcal{E}$}}] ++(0,{2*\h}) -- ++(4,0) to [resistor={info=$R_e$}] ++(0,-{2*\h}) -- (START);
		\node[contact] (1)  at (2.5,-\h) {}; \node[below] at (1) {$1$};
		\node[contact] (2)  at (2.5,+\h) {}; \node[above] at (2) {$2$};
		\draw[latex-latex, ultra thick, blue, font=\small] (1) -- node[left] {$U_{12}$} (2);
		\draw[-latex, ultra thick, red, font=\small, yshift=10pt] (0.5,-\h) -- node[above] {$I$} ++(1,0);
	\end{tikzpicture}
	\caption{Електричне коло}
	\label{thpic1}
\end{wrapfigure}
Електричний струм на ділянці виникає тоді, коли на її затискачах існує різниця потенціалів (електричне поле вздовж ділянки).

%\begin{Definition}
\emph{Напруга} (або падіння напруги) між двома точками~--- це різниця їхніх потенціалів. Для точок $1$ і $2$ (рис.~\ref{thpic1}): $U_{12} = \varphi_1 - \varphi_2$. Якщо $U_{12} > 0$, то $\varphi_1 > \varphi_2$, і струм тече від точки $1$ до точки $2$.

Напруга чисельно дорівнює роботі $A$, яку виконує джерело при переміщенні одиничного заряду між двома точками:
\begin{equation}\label{Udef}
	U_{12} = \frac{A_{12}}{q}.
\end{equation}
%\end{Definition}

Потужність~--- робота, яку виконує джерело за одиницю часу на даній ділянці кола:
\begin{equation}\label{Pdil}
	P = \frac{A_{12}}{t} = U_{12}I.
\end{equation}

%\begin{Definition}
\emph{Електрорушійна сила (ЕРС)} $\mathcal{E}$~--- це робота стороніх (неелектростатичних) сил джерела при переміщенні одиничного заряду по замкнутому колу:
\begin{equation}
    \mathcal{E} = \frac{A}{q}.
\end{equation}
Потужність, що розвивається джерелом:
\begin{equation}\label{Psource}
	P = \mathcal{E}I.
\end{equation}
%\end{Definition}

\begin{wrapfigure}{l}{0.4\linewidth}\centering
	\begin{tikzpicture}[thick, every circuit symbol/.style={thick},large circuit symbols,]
		\draw (0,0) node [contact] {} to [battery={info={$\mathcal{E}$}}] ++(2,0) to [resistor={info=$R_i$, color=red}] ++(2,0) node [contact] {} ;
	\end{tikzpicture}
	\caption{Еквівалентна схема джерела}
	\label{thpic2}
\end{wrapfigure}

%\begin{Definition}
\emph{Внутрішній опір} $R_i$~--- опір електричному струму всіх елементів всередині генератора (рис.~\ref{thpic2}).
%\end{Definition}

Якщо $R_i \ll R_e$, то напруга на клемах джерела практично не залежить від струму навантаження і дорівнює ЕРС. Таке джерело називають \emph{ідеальним джерелом напруги} ($R_i = 0$). Якщо ж $R_i \gg R_e$, то струм у зовнішньому колі практично не залежить від опору навантаження~--- це \emph{ідеальне джерело струму}. Надалі будемо вважати, що ЕРС та внутрішній опір джерела не залежать від струму.

В електричному колі (рис.~\ref{thpic1}) при зміні навантаження $R_e$ від нуля до нескінченності змінюються $I$ і $U$. Розглянемо найбільш характерні режими роботи джерела:

%\begin{ModeBox}
\begin{enumerate}
	\item \textbf{Режим холостого ходу:} $R_e \to \infty$, $I = 0$, $U = \mathcal{E}$.
	\item \textbf{Номінальний режим:} від джерела відбирається найбільша тривала потужність без перегрівання. Нагрівання зумовлено внутрішніми втратами: $P_i = I^2 R_i$.
	\item \textbf{Режим короткого замикання:} $R_e = 0$, $U = 0$, $I = I_\text{кз} = \mathcal{E}/R_i$. Для джерел з малим $R_i$ (акумулятори) це аварійний режим.
	\item \textbf{Режим узгодження:} $R_e = R_i$~--- максимальна потужність у навантаженні.
\end{enumerate}
%\end{ModeBox}

%\begin{LawBox}[title=ККД джерела]
Коефіцієнт корисної дії (ККД) $\eta$ визначається відношенням потужності в навантаженні ($P_e = I^2R_e$) до повної потужності джерела ($P = \mathcal{E}I$):
\begin{equation}\label{eta}
	\eta = \frac{P_e}{P} = \frac{I^2R_e}{I^2R_i + I^2R_e} = \frac{R_e}{R_i + R_e}.
\end{equation}
З формули видно: чим більше $R_e$ порівняно з $R_i$, тим вищий ККД.
%\end{LawBox}

%---------------------------------------------------------
\begin{wrapfigure}[19]{l}{0.45\linewidth}%[h!]
	\centering
	\begin{tikzpicture}[
			mark position/.style args={#1(#2)}{
					postaction={
							decorate,
							decoration={
									markings,
									mark=at position #1 with \coordinate (#2);
								}
						}
				}
		]
		\draw[-latex, thick] (0,0) -- ++(6.5,0) node[below] {$I$};
		\draw[-latex, thick] (0,0) -- ++(0,4) node[left] {$\eta$};
		\draw[-latex, thick] (6,0) node[below] {$I_\text{кз}$}-- ++(0,7) node[right] {$P$};
		\node[below] at (3,0) {$I_\text{кз}/2$};
		\node[left] at (0,3) {$1$};
		\draw[dashed] (3,0) -- (3,3/2);
		\draw[dashed] (0,3/2) node[left] {$0.5$} -- (3,3/2);
		\draw[ultra thick, domain=0:6, smooth, variable=\x, red,  mark position=0.8(A)]  plot ({\x}, {-1/6*(\x)*(\x-6)}) ;\node[above] at (A) {$P_e$};

		\draw[ultra thick, domain=0:6, smooth, variable=\x, blue,  mark position=0.2(E)] plot ({\x}, {-1/2*\x+3}); \node[above] at (E) {$\eta$};

		\draw[ultra thick, domain=0:6, smooth, variable=\x, green!60!black,  mark position=0.7(Pe)] plot ({\x}, {1/6*\x*\x}); \node[anchor=north west] at (Pe) {$I^2R_i$};

		\draw[ultra thick, domain=0:6, smooth, variable=\x, green!60!black,  mark position=0.7(P)] plot ({\x}, {1*\x});\node[anchor=south east] at (P) {$P$};
	\end{tikzpicture}
	\caption{Діаграма режимів роботи джерела електроенергії}
	\label{thpic3}
\end{wrapfigure}
%---------------------------------------------------------
При узгодженому режимі ($R_e = R_i$) у приймачі виділяється найбільша потужність. Цей режим використовується у вимірювальних колах, пристроях обчислювальної та інформаційної техніки, засобів зв'язку. В цьому випадку $\eta = 1 - \frac{I}{I_\text{кз}}$, де $I_\text{кз}$~--- струм короткого замикання. Потужність у навантаженні при цьому:
\begin{equation}\label{key}
	P_e = \frac{\mathcal{E}^2}{4R_i},
\end{equation}
а ККД дорівнює $50$~\%.

%\begin{NoteBox}
При передачі великих потужностей (енергетика) режим узгодження є неприпустимим~--- там ККД має бути якнайвищим. Обов'язкова умова: $R_e \gg R_i$.
%\end{NoteBox}

На рис.~\ref{thpic3} наведено залежності $\eta$, $P_e$, $I^2R_i$ та $P$ від струму $I$. З кривих видно, що найбільша потужність у навантаженні досягається при $I = I_\text{кз}/2$, або що те ж~--- при $R_e = R_i$.

\section{Теоретичне підґрунтя роботи}

\subsection{Закони Ома та ідея дослідів}

%---------------------------------------------------------
\begin{wrapfigure}[12]{l}{0.5\linewidth}\centering
	\begin{tikzpicture}[thick, every circuit symbol/.style={thick},large circuit symbols,]
		\draw[red, fill=red!20] (-0.4,-1.4) rectangle ++(0.8,2.8);
		\draw (0,-1.5) coordinate (START) to [resistor={info={[label distance=1mm]$R_i$}}] ++(0,1.5) to [battery={info={$\mathcal{E}$}}] ++(0,1.5) -- ++(1,0) to [ampermeter] coordinate (A1) ++(3,0) to [resistor={info'=$R_e$}] ++(0,-3) -- (START)
		;
		\draw [latex-] (4.1,0) -- ++(0.5,0) -- ++(0,-1) -- ++(-0.6,0) node[contact] {};
		\draw ([xshift=1cm]START) node[contact] {} to[voltmeter] ++(0,3) node[contact] {};
	\end{tikzpicture}
	\caption{Послідовне з'єднання елементів}
	\label{pic1}
\end{wrapfigure}
%---------------------------------------------------------
Реальні джерела описуються в еквівалентній схемі як послідовно з'єднані ідеальне \emph{джерело напруги} $\mathcal{E}$ і внутрішній опір $R_i$ (рис.~\ref{pic1}). Напруга на клемах і сила струму залежать від зовнішнього опору $R_e$.

%\begin{LawBox}[title=Закон Ома для повного кола]
\begin{equation}\label{OhmLow}
	I = \frac{\mathcal{E}}{R_i + R_e}.
\end{equation}


Вихідна напруга (напруга на клемах джерела):
\begin{equation}\label{U}
	U = \mathcal{E} - R_i I.
\end{equation}


Зовнішній опір у робочій точці визначається як:
\begin{equation}\label{Re}
	R_e = \frac{U}{I}.
\end{equation}
%\end{LawBox}


Формула~\eqref{U} описує лінійну залежність $U(I)$, де вільний член дає $\mathcal{E}$ (ЕРС~--- напруга при $I=0$, тобто в режимі холостого ходу), а нахил прямої дорівнює $-R_i$. Тому:
\begin{equation}\label{tan}
	\tg\varphi = - R_i,
\end{equation}
де $\varphi$~--- кут нахилу залежності $U(I)$ до осі $I$ (рис.~\ref{pic2}).

\begin{NoteBox}
Не плутати два різні нахили на графіку $U(I)$:\par
\textbf{1)} Нахил самої прямої $U(I)$~--- визначає $R_i$ через $\tg\varphi = -R_i$.\par
\textbf{2)} Нахил прямої від початку координат до робочої точки~--- визначає зовнішній опір $R_e = U/I$ в цій точці.
\end{NoteBox}

%---------------------------------------------------------
\begin{figure}[h!]\centering
	\begin{tikzpicture}
		\draw[-latex] (0,0) -- ++(5,0) node[below] {$I$, А};
		\draw[-latex] (0,0) -- ++(0,5) node[left] {$U$, В};
		\draw[red, thick] (0,4) -- (4,0) coordinate[pos=0.5] (X) coordinate[pos=0.4] (A);
		\draw[dashed] (X) -- (0,0-|X) node[below] {$I_\text{кз}/2$};
		\draw[dashed] (X) -- (0,0|-X) node[left] {$\mathcal{E}/2$};
		\fill (X) circle (0.1);
		\node[right] at (X) {$R_e = R_i$};
		\node[left] at (0,4) {$\mathcal{E}$};
		\node[below] at (4,0) {$I_\text{кз}$};
		\node[right] at (0,4) {$R_e \to \infty$};\node[above right] at (4,0) {$R_e \to 0$};
		\draw (3,0) arc(180:{180-45}:1) node[anchor=north east, xshift=-5pt, font=\small] {$\pi - \varphi$};

        \draw (0,0) -- (X);
	\end{tikzpicture}
	\caption{Характеристика джерела із сталим внутрішнім опором $R_i$.}
	\label{pic2}
\end{figure}
%---------------------------------------------------------

Таким чином, лінійна залежність $U(I)$ означає, що внутрішній опір $R_i$ є сталим.

Для нашої роботи ми будемо аналізувати три типи навантаження:

\begin{enumerate}
	\item \textbf{Холостий хід:} $R_e \to \infty$. Струм не тече, немає падіння напруги на $R_i$, тому $U = \mathcal{E}$.
	\item \textbf{Коротке замикання:} $R_e=0$. Все падіння напруги~--- на внутрішньому опорі, тому $U = 0$. Струм короткого замикання:
	      \begin{equation}\label{short}
		      I_\text{кз}=\frac{\mathcal{E}}{R_i}.
	      \end{equation}
	\item \textbf{Узгодження потужності:} $R_e=R_i$. У цьому випадку $U = \frac{\mathcal{E}}{2}$, $I=\frac{I_\text{кз}}{2}$.
\end{enumerate}

Пронормуємо виміряні значення напруги $U$ та струму $I$ на максимальні значення~--- напругу холостого ходу $\mathcal{E}$ та струм короткого замикання $I_\text{кз}$. З рівнянь \eqref{OhmLow} та \eqref{short}:
\begin{equation}\label{I/I0}
	\frac{I}{I_\text{кз}} = \frac{1}{1 + \dfrac{R_e}{R_i}}.
\end{equation}

Рівняння \eqref{U} дає
\begin{equation}\label{Г/Г0}
	\frac{U}{\mathcal{E}} = 1 - \frac{1}{1 + \dfrac{R_e}{R_i}},
\end{equation}
або, враховуючи \eqref{I/I0}:
\begin{equation}
	\frac{U}{\mathcal{E}} + \frac{I}{I_\text{кз}}  = 1.
\end{equation}

Тобто нормовані залежності напруги та струму від навантаження є взаємно протилежними (рис.~\ref{pic3}).

\subsection{Потужність та ККД}

Визначимо співвідношення між потужностями елементів кола.

Максимальна потужність кола (при короткому замиканні):
\begin{equation}\label{eq:P0}
    P_0 = \mathcal{E} I_\text{кз}.
\end{equation}

Потужність на внутрішньому опорі $R_i$:
\begin{equation}\label{Pi}
	P_i = I^2R_i.
\end{equation}

Потужність на навантаженні $R_e$:
\begin{equation}\label{Pe}
	P_e = I^2R_e.
\end{equation}

Нормуємо потужність навантаження~\eqref{Pe} на максимальну~\eqref{eq:P0}:
\begin{equation}\label{Pe/P0}
	\frac{P_e}{P_0} = \frac{U}{\mathcal{E}} \cdot \frac{I}{I_\text{кз}} = \frac{R_e/R_i}{\left( 1+ R_e/R_i \right)^2 }.
\end{equation}

Залежність $P_e/P_0$ від $R_e/R_i$ має максимум при $R_e=R_i$ (рис.~\ref{pic3})~--- джерело віддає в навантаження максимальну потужність у режимі узгодження.

%---------------------------------------------------------
\begin{figure}[h!]\centering
	\begin{tikzpicture}
		\begin{semilogxaxis}[
				x label style={at={(axis description cs:0.5,-0.1)},anchor=north},
				legend style={row sep = 5pt},
				xlabel={$\frac{R_e}{R_i}$},
				width = \linewidth,
				height = 0.6\linewidth,
				grid = both,
				major grid style={line width=.6pt,draw=brown!60},
				minor tick num = 9,
				minor grid style = {line width=.1pt,draw=brown!20},
				axis lines=center,
			]
			\addplot[domain=0.01:100, samples=100, blue, thick] {1 - 1/(1 + x/1)};
			\addplot[domain=0.01:100, samples=100, green!50! black, thick] {1/(1 + x/1)};
			\addplot[domain=0.01:100, samples=100, red, thick] {x/1/(1+x/1)^2};
			\node (XX) [draw=red, fill= white, minimum width=5cm] at (axis cs:10,0.5) {Холостий хід};
			\node (KZ) [draw=red, fill= white, minimum width=5cm] at (axis cs:0.1,0.5) {Коротке замикання};
			\draw[-latex, double] (XX.east) -- ++(0.5,0);
			\draw[-latex, double] (KZ.west) -- ++(-0.5,0) ;
			\legend{${U}/{\mathcal{E}}$, ${I}/{I_\text{кз}}$, ${P}/{P_0}$}
		\end{semilogxaxis}
	\end{tikzpicture}
	\caption{Діаграма режимів для джерела електроенергії}
	\label{pic3}
\end{figure}
%---------------------------------------------------------

\section{Експериментальне устаткування}

Експериментальне устаткування складається з трьох типів джерел електричної енергії:
\begin{itemize}
    \item \textbf{свинцево-кислотний акумулятор} (~$12$~В)~--- джерело з малим внутрішнім опором;
    \item \textbf{плоска батарея} $4{,}5$~В (три послідовно з'єднані елементи Лекланше)~--- джерело із середнім внутрішнім опором;
    \item \textbf{нікель-кадмієвий акумулятор} ($1{,}2$~В)~--- джерело з дуже малим внутрішнім опором.
\end{itemize}
Як змінне навантаження використовується реостат на $100$~Ом або на $10$~Ом (залежно від типу джерела). Вимірювальні прилади~--- високоомний цифровий мультиметр (як вольтметр) та магнітоелектричний мультиметр класу точності $0{,}5$ (як амперметр).

\section{Хід експерименту}

\subsection*{Частина~1. Зняття вольт-амперної характеристики}

\begin{enumerate}
    \item Зберіть вимірювальну схему згідно з рис.~\ref{pic1}. Як змінний резистор $R_e$ використовуйте реостат на $100$~Ом.

    \item Встановіть реостат у положення максимального опору ($R_e \to \infty$) і виміряйте напругу холостого ходу $U_0 \approx \mathcal{E}$.

    \item Змінюючи опір реостата, знімайте залежність $U(I)$~--- від режиму холостого ходу ($I = 0$) до максимально допустимого струму з кроком $\Delta I \approx 0{,}1$~А.

    \item Занесіть результати до таблиці~\ref{tab1}. Для кожної точки обчисліть $R_e = U/I$.
\end{enumerate}

\begin{center}
    \captionof{table}{Вимірювання вольт-амперної характеристики}
    \label{tab1}
    \begin{tblr}{
    colspec={Q[c, 1.45cm] Q[c, 1.45cm] Q[c, 1.46cm] Q[c, 1.45cm] Q[c, 1.45cm] Q[c, 1.45cm] Q[c, 1.45cm] Q[c, 1.45cm] Q[c, 1.45cm]},
    hlines, vlines,
    }
        $I$, А & $U$, В & $R_e$, Ом & $R_i$, Ом & $R_e/R_i$ & $P_e/P_0$ & $P_i / P_0$ & $P/ P_0$ & $\eta$ \\
              &       &       &     &           &           &             &          &
    \end{tblr}
\end{center}

\subsection*{Частина~2. Визначення параметрів джерела}

\begin{enumerate}
    \item Побудуйте графік залежності $U(I)$.

    \item Апроксимуйте отримані точки прямою лінією $U = \mathcal{E} - R_i I$. \emph{Підказка:} ЕРС $\mathcal{E}$~--- це значення $U$ при $I = 0$ (точка перетину з віссю $U$); внутрішній опір $R_i$~--- це модуль нахилу прямої:
    \begin{equation*}
        R_i = -\frac{\Delta U}{\Delta I}.
    \end{equation*}

    \item Визначте струм короткого замикання $I_\text{кз} = \mathcal{E}/R_i$ (точка перетину прямої з віссю $I$).

    \item Позначте на графіку точку узгодження ($R_e = R_i$, тобто $I = I_\text{кз}/2$, $U = \mathcal{E}/2$).
\end{enumerate}

\subsection*{Частина~3. Аналіз потужностей та ККД}

\begin{enumerate}
    \item Обчисліть $P_0 = \mathcal{E} \cdot I_\text{кз}$ (максимальна потужність джерела).

    \item Для кожної точки таблиці обчисліть нормовані величини $P_e/P_0$, $P_i/P_0$, $P/P_0$ та ККД $\eta$ і занесіть до таблиці~\ref{tab1}.

    \item Побудуйте нормовані діаграми потужностей у функції від $R_e/R_i$, аналогічні рис.~\ref{pic3}. Використайте логарифмічну шкалу по осі $x$.

    \item Порівняйте вигляд отриманого графіка з теоретичним (рис.~\ref{pic3}) та прокоментуйте розбіжності.
\end{enumerate}

\subsection*{Частина~4. Класифікація джерела}

\begin{enumerate}
    \item Порівняйте $R_i$, отриманий з апроксимації, з типовими значеннями $R_e$ з таблиці вимірювань: чи $R_i \ll R_e$, чи $R_i \sim R_e$, чи $R_i \gg R_e$?
    \item Зробіть висновок: чи наближається досліджуване джерело до ідеального джерела напруги ($R_i \ll R_e$), чи до ідеального джерела струму ($R_i \gg R_e$)?
\end{enumerate}


\section*{Контрольні питання}

\begin{enumerate}
    \item Що таке електрорушійна сила (ЕРС) джерела та який її фізичний зміст?

    \item Чим відрізняється ЕРС від напруги на клемах джерела?

    \item Запишіть закон Ома для повного кола.

    \item Що таке внутрішній опір джерела і як він впливає на струм у колі?

    \item Як змінюється напруга на клемах джерела при збільшенні струму?

    \item Що таке режим холостого ходу і які значення струму та напруги при цьому?

    \item Що таке коротке замикання? Запишіть формулу струму короткого замикання.

    \item Як експериментально визначити ЕРС і внутрішній опір джерела?

    \item Який вигляд має графік залежності напруги від струму $U(I)$ і що означає його нахил?

    \item За якої умови потужність, що виділяється на навантаженні, є максимальною?
\end{enumerate}
